\chapter{Palpitations}

\section*{Definition}
Awareness of heartbeat perceived as rapid, pounding, fluttering, or skipped beats.

\section*{What Happens in Your Body}
Often benign but can indicate arrhythmias, structural heart disease, metabolic abnormalities, anxiety, or stimulant effects.

\section*{Causes}
\begin{itemize}
  \item Anxiety/stress
  \item Arrhythmias (PVCs, SVT)
  \item Electrolyte imbalance
  \item Hyperthyroidism
  \item Caffeine/medication
\end{itemize}

\section*{When to See a Doctor}
See your doctor if:
\begin{itemize}
  \item Symptoms are severe or getting worse over time
  \item Evaluate if new, irregular, or accompanied by chest — seek medical evaluation
  \item You have other concerning symptoms such as fever, weight loss, or bleeding
\end{itemize}

\section*{Related Symptoms}
\begin{itemize}
  \item Dizziness
  \item Chest pain
  \item Shortness of breath
  \item Sweating
  \item Syncope/fainting
\end{itemize}
\textbf{Important:} If this symptom persists, your doctor will check for serious underlying causes.

\section*{Self-Care / Home Management}
\begin{itemize}
  \item Monitor symptoms including frequency, duration, triggers, and associated features; keep a symptom log.
  \item Avoid stimulants (caffeine, nicotine, decongestants, energy drinks) which can exacerbate palpitations and tachycardia.
  \item Practice relaxation techniques and stress reduction to minimize sympathetic activation.
  \item Maintain adequate hydration and electrolyte balance, especially during illness or exercise.
  \item Seek immediate evaluation for chest pain, shortness of breath, fainting, or palpitations with dizziness.
\end{itemize}

\section*{How Your Doctor Will Evaluate You}
\begin{itemize}
  \item History characterizes the symptom (quality, onset, duration, triggers, relieving factors) and assesses cardiac risk factors.
  \item Physical examination includes vital signs, cardiac auscultation (rate, rhythm, murmurs), lung exam, and peripheral edema assessment.
  \item Diagnostic testing includes ECG, ambulatory cardiac monitoring (Holter or event monitor), echocardiogram, and stress testing.
  \item Laboratory studies (troponin, BNP, electrolytes, thyroid function) help identify acute or contributing pathology.
\end{itemize}

\section*{Treatment Options}
\begin{itemize}
  \item Benign palpitations require reassurance and trigger avoidance; persistent symptoms may benefit from beta-blockers.
  \item Arrhythmias are treated based on type: rate control, rhythm control, anticoagulation, or ablation procedures.
  \item Heart murmurs require echocardiographic evaluation; management depends on the specific valvular lesion and severity.
  \item Cardiology referral is indicated for concerning symptoms, abnormal ECG findings, or newly discovered murmurs.
\end{itemize}

\section*{References}
\begin{thebibliography}{99}
\bibitem{palpitations:palp1} Zimetbaum P, Josephson ME. ``Evaluation of patients with palpitations.'' N Engl J Med. 1998;338(19):1369-1373.
\bibitem{palpitations:palp2} Thavendiranathan P, Bagai A, Khoo C, et al. ``Does this patient with palpitations have a cardiac arrhythmia?'' JAMA. 2009;302(19):2135-2143.
\bibitem{palpitations:palp3} Wexler RK, Pleister A, Raman SV. ``Outpatient approach to palpitations.'' Am Fam Physician. 2011;84(1):63-69.
\bibitem{palpitations:palp4} Raviele A, Giada F, Bergfeldt L, et al. ``Management of patients with palpitations: a position paper from the EHRA.'' Europace. 2011;13(7):920-934.
\bibitem{palpitations:palp5} Zimetbaum PJ, Josephson ME. ``Approach to the patient with palpitations.'' UpToDate, 2023.
\bibitem{palpitations:palp6} Giada F, Raviele A. ``Clinical approach to patients with palpitations.'' Card Electrophysiol Clin. 2018;10(2):303-312.
\end{thebibliography}
